\documentclass[12pt]{article}

\usepackage[T1]{fontenc}
\usepackage{lmodern}
\usepackage{setspace}
\usepackage{geometry}
\usepackage{microtype}
\usepackage[hidelinks]{hyperref}
\usepackage{parskip}
\usepackage{csquotes}

\usepackage[
  backend=biber,
  style=authoryear,
  citestyle=authoryear
]{biblatex}

\addbibresource{thesis_bibliography_short.bib}

\geometry{margin=1in}
\onehalfspacing

\title{The Invisible Limits of Public Reason:\\
Stability, Tracking, and the Political Epistemology of Justification\\
\normalsize A thesis submitted in partial fulfillment of the requirements for the degree of Master of Arts in Philosophy\\
\normalsize Department of Philosophy\\
\normalsize San Francisco State University\\}
\author{Alex Horovitz}
\date{March 2026}

\begin{document}
\maketitle

\section{Introduction: The Puzzle}

For a time it became fashionable to speak as if liberal democracy had solved its legitimacy problem. That confidence now reads like a period piece. A liberal public order can fail without a coup. It can fail because citizens no longer experience coercive law as something they can endorse as free and equal co-authors. That is why stability has returned as a central political problem in most people's living memory \parencite[p.~1187]{Weithman2024Rescuing}.

Rawlsian political liberalism remains the canonical attempt to explain how legitimacy is possible in a society that is both free and pluralistic. Rawls begins from a fact about modern liberal life that he treats as permanent rather than regrettable: people will affirm incompatible comprehensive doctrines under conditions of liberty. The project is to show how such citizens can nevertheless endorse a shared political conception of justice, and to show that this endorsement can be stable for the right reasons rather than a mere truce \parencite[p.~1]{rawls1987overlapping}. A modus vivendi can be stable through fear or exhaustion. Rawls wants a stability that expresses mutual respect.

This thesis offers an internal critique of that ambition. I am not arguing that political philosophy should be replaced by psychology, nor that norms of public justification are illegitimate because citizens are biased or selfish. Rawls's legitimacy route assigns citizens an epistemic role: they must be able to recognize the justificatory status of coercive laws through reasons they can regard as public. Once legitimacy depends on that role, feasibility conditions appear from inside the normative project. They appear because the theory requires citizens to be able to do something, not because I have imported a hostile realist standpoint.

The guiding question is therefore simple. What must be true about citizens, institutions, and political environments for stability through public justification to work as Rawls needs it to work? This question has a second edge. If those conditions do not obtain, what follows for a Rawlsian account of legitimacy?

I defend two conclusions that must be kept distinct. The first is a feasibility claim: under modern conditions of mediation and rapid drift, there may be no institutional form that reliably delivers Rawlsian stability-through-shared-public-justification at scale. The second is an internal-critique claim: if Rawlsian legitimacy is defined by stability for the right reasons, then limits on citizens' ability to play the required epistemic role are limits internal to the ideal. These are not external debunking arguments. They are constraints on the kind of justificatory practice the ideal requires.

Two clarifications are needed before the machinery begins. First, the argument is not that political life is always too turbulent for any stable justification. The burden changes depending on the time horizon over which stability is demanded. A claim about short-cycle drift is not automatically a claim about long-cycle stability. Second, the argument is not that public reason requires citizens to evaluate everything at full policy resolution. The strongest Rawlsian reply is that public reason operates at the level of constitutional essentials. I accept that reply and treat it as a constraint on my own critique. The question is what still survives after the demand is made coarse enough to be faithful to Rawls's own restrictions.

A short vignette sharpens the stakes. Consider the justificatory terrain around privacy, speech, and security. The moral language stays familiar: liberty, equal citizenship, the public good, the limits of state power. Yet the empirical predicates shift quickly. Changes in surveillance technology, changes in the perceived threat environment, and shifts in what counts as a reasonable inference from data can alter the practical meaning of policies that keep their old names. A law that once looked like a narrow investigative tool can become, under new technical conditions, a permanent infrastructure of monitoring. Citizens can remain sincere and normatively reasonable while failing to track that transformation in time. The claim is not that citizens refuse to deliberate. The claim is that the world can change in ways that defeat the epistemic work that stability through justification requires them to do.

The argument proceeds in four stages. I reconstruct Rawlsian stability and public reason as an epistemic assignment to citizens and officials. I then make two implicit assumptions explicit: citizens must be able to track justificatory status across time, and the justificatory landscape must be stable enough for that tracking to work. I show why those assumptions are fragile in mediated politics, drawing on work on information cascades, polarization dynamics, bounded cognition, and social complexity. I end by arguing that if a Rawlsian spirit of public justification is to survive under drift, it will require institutional scaffolding that reduces citizens' tracking burden without replacing their judgment. That constructive proposal invites a paternalism objection, and the thesis must earn its answer.

\section{Rawlsian Stability and Public Reason}

Rawls binds legitimacy to a justificatory discipline. The state may coerce, but only under terms that free and equal citizens can accept as reasonable. In \emph{Public Reason Revisited} Rawls presents public reason as a political ideal that governs how fundamental political matters should be decided \parencite[p.~765]{rawls1997public}. It is an ideal about the way a constitutional democracy should govern itself, not simply about which outcomes are morally attractive.

Rawls's own formulation is explicitly interpersonal. In exercising political power, we should not merely have reasons; we must be able to offer reasons that satisfy a reciprocity constraint. Rawls writes that our exercise of political power is proper only when we sincerely believe the reasons we would offer for political action are sufficient and we reasonably think other citizens might also reasonably accept those reasons \parencite[pp.~446--447]{rawls2005political}. Reciprocity is not a demand for agreement. It is a demand that coercion be offered under terms that address others as moral equals.

That demand reaches beyond officials. Rawls connects public reason to a duty of civility that shapes political conduct and self-understanding. In a representative democracy citizens exercise political power by voting, and Rawls says they should think of themselves as ideal legislators, asking what statutes and policies can be supported by reasons that satisfy reciprocity. They should also be disposed to repudiate officials who violate public reason \parencite[p.~444]{rawls2005political}. The point is not that every citizen must become a constitutional lawyer. The point is that citizens do not merely receive justification. They participate in sustaining it.

Rawls insists that public reason has a definite subject matter. Political conceptions of justice should be complete in a practical sense: they should specify values and principles that can answer \enquote{all or nearly all} questions involving constitutional essentials and matters of basic justice \parencite[pp.~453--454]{rawls2005political}. This narrowing is often invoked as a reply to feasibility worries. Public reason is not supposed to govern every policy detail, so citizens are not asked to track everything.

Stability for the right reasons is what gives this architecture its bite. Rawls does not want a society that remains orderly only because citizens fear punishment, or because the balance of power makes defection too costly. He wants a society whose basic structure can be justified to citizens in a way that connects their compliance to reasons they can endorse from within their own doctrines. That is why \emph{overlapping consensus} matters. It is not agreement on a comprehensive worldview. It is the idea that citizens with different doctrines can affirm the same political conception of justice, each from within their own doctrine \parencite[p.~8]{rawls1987overlapping}. It is also meant to secure a distinctive kind of stability. Citizens should not merely acquiesce in a constitutional order. They should be able to explain to themselves why its coercion is justified, and they should be able to recognize that their fellow citizens can make the same explanation in their own terms.

Rawls also connects public reason to civic affect. He worries that when citizens no longer see the point of affirming the ideal of public reason, they may fall into bitterness and resentment \parencite[p.~485]{rawls2005political}. That is a striking remark in context. It treats the ideal not as a sterile rule for argument, but as part of the social psychology of a constitutional democracy. If citizens cannot sustain the epistemic practice that public reason requires, the failure is not merely theoretical. It shows up as an erosion of the moral sentiments that make stability possible.

The burdens of judgment are meant to complete the picture. They explain why reasonable people will disagree under conditions of liberty, without treating disagreement as a sign of irrationality or bad faith. The burdens are crucial for Rawls's moral aim: citizens can respect opponents as reasonable even when they believe those opponents are mistaken about the good. Quong defends this structure against a common critique, arguing that political liberalism's demand for epistemic restraint does not assume skepticism about the good. It is motivated by a moral desire to justify coercion to others \parencite[p.~322]{quong2007without}. The philosophical point is subtle. Public reason does not ask citizens to doubt their faith or their moral commitments. It asks them to recognize what they owe to others as co-citizens in a shared coercive order.

That brings me to a distinction I will rely on throughout. Rawls's \emph{reasonable} is not a psychological type. It is a normative stance: a willingness to propose fair terms of cooperation and to accept the burdens of judgment as a constraint on what can be imposed through collective power. The \emph{rational} is also a normative construct in Rawls, capturing the pursuit of one's ends within fair terms. My critique does not rest on a claim that citizens are unreasonable. One can accept Rawls's moral psychology at the level of motivation and still doubt that the epistemic tasks embedded in stability-through-justification can be carried out under real political conditions.

A second point matters for later chapters. Rawlsian stability is often treated as static: a regime is stable if citizens endorse it. Weithman presses a more sophisticated picture. Stability is an equilibrium notion. It concerns what citizens can reliably do over time given their beliefs, motivations, and expectations about one another. When Rawls concedes that an overlapping consensus may emerge over a family of liberal conceptions rather than a single uniquely correct one, the stability problem changes shape. Stability becomes set-valued rather than point-valued: a society can be well ordered so long as it remains within a range of acceptable liberal conceptions, not because it has settled on a single fixed point \parencite[pp.~205--206]{Weithman2023FixedPoints}. This shift will later become the strongest Rawlsian reply to my critique. It deserves to be presented with full force.

\section{Feasibility as an Internal Constraint}

Feasibility objections in political philosophy usually come in familiar forms. Some concern motivation: people will not comply even if they know what justice requires. Some concern institutions: power will be captured, rules will be enforced unevenly, incentives will corrupt. Some concern information: citizens will be ignorant, elites will manipulate. Others concern dynamics: even if a design works at one time, it may not remain stable as circumstances change. These objections are often treated as external constraints on an otherwise intact ideal.

That framing is not quite right for Rawlsian legitimacy. Rawls is not offering a purely aspirational ideal of justice detached from social life. He treats stability for the right reasons as part of what makes a political conception morally compelling. If a conception cannot plausibly generate its own support under conditions it claims to respect, it has a justificatory defect, not merely an implementation problem. This is why feasibility enters from within. If legitimacy depends on citizens being able to endorse coercive law for public reasons, then whatever undermines that ability undermines legitimacy, not merely compliance.

The point is easiest to see by contrast. Some ideal theorists deny that feasibility constraints matter for the truth of principles. Cohen defends the thought that ultimate principles can be fact-insensitive even if rules of regulation must accommodate facts \parencite[p.~20]{cohen2008rescuing}. Estlund resists what he calls \enquote{utopophobia} by insisting that political philosophy can articulate ideals even when they are not fully realizable \parencite[p.~4]{estlund2019utopophobia}. Valentini reconstructs the debate as a paradox: ideal theory seems indispensable for guidance, yet critics worry that ideal assumptions make it unable to guide action \parencite[p.~332]{valentini2009paradox}. These debates concern the standing of normative principles in general.

Rawlsian legitimacy is a narrower target. It is not merely a claim about what justice requires. It is a claim about what makes coercion legitimate for citizens who disagree. That legitimacy claim is defined by a stability-through-justification ambition. Once one ties legitimacy to that ambition, one cannot treat feasibility as an external embarrassment. The theory's success depends on citizens and institutions being able to do specific things.

Publicity makes the internality vivid. Freeman argues that without publicity moral principles cannot serve as principles of practical reasoning and justification among free and equal persons, and that without a public basis of justification political discussion becomes susceptible to manipulation by persuasive rhetoric and familiar biases \parencite[p.~11]{freeman2007burdens}. Freeman is not offering my argument, but he captures the same dependence: justification is supposed to be public in a way that citizens can use to orient themselves, criticize, and demand reasons.

Two temptations arise at this point. One is to retreat into a purely institutional story: perhaps courts, legislatures, and public norms can carry the justificatory burden, so individual limits do not matter. The other is to retreat into an epistemic elitism: perhaps citizens can defer to experts who know what public reasons support. Both temptations are unstable.

The institutional story cannot fully bypass citizens. Rawls ties public reason to civic self-understanding and to the duty of civility. Citizens are asked to think of themselves as participants in public justification, not as subjects who obey because institutions have spoken \parencite[p.~444]{rawls2005political}. A legitimacy route that turns citizens into passive recipients ceases to be Rawlsian in the relevant sense.

Expert deference does not solve the problem either. Tetlock's study of expert political judgment finds few signs that expertise reliably translates into better calibrated forecasting in the domains where experts are most confident \parencite[p.~20]{tetlock2005expert}. The point is not that expertise is worthless. It is that the dream of replacing a citizen-facing justificatory practice with a technocratic surrogate runs into limits of its own. If expertise is unreliable, citizens cannot simply outsource the tracking problem.

A third route aims to reduce the demand by loosening the idea of shared public reasons. Vallier's convergence liberalism is the most systematic version of this move. Vallier proposes combining a convergence conception of public reasons with a moderate conception of idealization, arguing that reasons offered to justify coercion need not be shared \parencite[p.~8]{vallier2011diss}. This does lower a familiar burden. Citizens need not filter themselves down to one shared justificatory currency. Yet it does not eliminate the need for tracking. If citizens can justify the same law from different reasons, they still need to track whether the law remains justified over time, and they still need public ways of identifying whether the justificatory structure has drifted.

The internal-critique claim I defend is therefore narrow. I am not denying that ideals can outrun feasibility. I am arguing that Rawls's particular legitimacy route makes feasibility internal because stability is doing justificatory work. Once one sees that, two implicit assumptions come into view.

\section{The Two Invisible Assumptions}

The first assumption is what I call the tracking requirement. It is a diachronic epistemic task. Citizens must be able to tell, as conditions change, whether coercive laws remain publicly justifiable. They need not re-derive the principles of justice from scratch each time a policy dispute arises. The demand is weaker but still substantial. To experience coercion as legitimate under Rawls's ideal, citizens must have some stable way of identifying whether the basic structure remains within the bounds of a liberal political conception and whether officials are acting on public reasons rather than sectarian ones. Rawls's description of public reason as a public political ideal suggests this kind of citizen-facing uptake \parencite[p.~767]{rawls1997public}.

Tracking is not simply knowing what the law says. It is knowing what the law does in relation to public political values, and recognizing when changes in empirical conditions alter that relation. It is also knowing when one has reason to treat a contested policy as still inside the liberal family rather than as a departure from it. Weithman's equilibrium framing makes clear why this is not optional. If stability is sustained by citizens' beliefs and expectations, citizens must be able to update those beliefs and expectations in response to institutional change, otherwise the stability story becomes a stipulation \parencite[p.~207]{Weithman2023FixedPoints}.

The second assumption is landscape stability. Even if citizens are willing to track, tracking can only succeed if the mapping from reasons and facts to policy consequences and justificatory status is stable enough over time. The thesis's metaphor of a justificatory landscape is meant to capture that mapping. A stable landscape is one in which the same public reasons reliably point toward roughly the same kinds of outcomes, and in which shifts in empirical predicates are slow enough, and legible enough, that citizens can adjust their judgments without repeated breakdowns.

Two parameters matter here, and they matter because the problem can be overstated if one ignores them. The first is timescale. The burden stability places on citizens looks different depending on whether we are asking for stability over a few years or over a generation. Short-cycle drift includes agenda shocks, crises, rapid technological changes, and fast institutional reinterpretations. Long-cycle drift includes generational replacement, slow evolution of doctrines, and long-run institutional path dependence. Weithman notes that stability worries have returned with force in part because these timescales can interact: short-cycle shocks can reset long-cycle expectations \parencite[p.~1188]{Weithman2024Rescuing}.

The second parameter is granularity. The strongest Rawlsian reply is that public reason operates at a coarse level. Citizens are not required to track every policy detail. They need only track whether constitutional essentials and basic justice remain within a liberal family of acceptable conceptions. This is not a weak objection. It is a plausible reconstruction of Rawls's own restriction of public reason's subject matter \parencite[pp.~453--454]{rawls2005political}. It is also what motivates Weithman's set-valued stability account, which aims to make this coarse-grained picture philosophically coherent \parencite[p.~205]{Weithman2023FixedPoints}.

My claim is not that the granularity reply fails. My claim is that it relocates the burden. Even if one grants that citizens need only track coarse-grained liberal family membership, they still need reliable public markers of what counts as inside that family. They also need a sufficiently stable mapping from constitutional values to concrete institutional outcomes for the family-membership judgment to be more than a slogan. Coarse-graining does not eliminate landscape instability. It changes its form.

\section{Why the Assumptions Fail in Mediated Politics}

A liberal society can fail the tracking requirement without citizens becoming unreasonable. The failure can arise because the epistemic environment disrupts the connection between public reasons and the public's grasp of what those reasons mean in practice. Two mechanisms matter.

The first is the shifting of empirical predicates through mediated diffusion. Political beliefs often move by cascades rather than by individual evaluation of evidence. Bikhchandani, Hirshleifer, and Welch describe how informational cascades can lead individuals to follow earlier movers even when their private signals diverge, because the social signal becomes the dominant input \parencite[p.~992]{bikhchandaniHirshleiferWelch1992cascades}. In a mediated public sphere the relevant early movers are not only neighbors but influencers, partisan media, and political entrepreneurs. Once a cascade forms, the public meaning of a policy label can shift quickly, and citizens can be drawn into patterns of belief that are locally rational relative to the social signal but poorly connected to the underlying facts.

Bikhchandani and colleagues emphasize that cascades can begin from small informational asymmetries and then become self-reinforcing because later agents rationally treat earlier actions as evidence \parencite[p.~996]{bikhchandaniHirshleiferWelch1992cascades}. That matters for public justification because it can change which empirical claims are treated as common knowledge. Once a certain description of a policy becomes the salient public description, it becomes harder for citizens to separate the justificatory structure from the rhetoric that now rides on top of it.

Online diffusion intensifies this problem. Vosoughi, Roy, and Aral find that false news on Twitter diffuses farther, faster, deeper, and more broadly than true news, with effects especially pronounced for political news \parencite[p.~1146]{vosoughiRoyAral2018falseNews}. They also report an affective difference: false stories inspire emotions like fear, disgust, and surprise, while true stories inspire trust and anticipation \parencite[p.~1146]{vosoughiRoyAral2018falseNews}. The point for my argument is not a moral panic about social media. It is that the empirical predicates citizens would need to track the justificatory status of policies are not presented in a neutral stream. They are filtered through mechanisms that reward novelty and emotional intensity. When the empirical substrate shifts through these mechanisms, citizens can endorse and defend a justificatory picture that has quietly detached from the policy reality they face.

The second mechanism is the shifting of moral salience through polarization dynamics. Levendusky and Malhotra show that media coverage of polarization increases citizens' beliefs that the electorate is polarized, and that this has a complex effect. It can moderate issue positions while increasing dislike of the opposing party \parencite[pp.~283--284, 291--293]{levendusky2016media}. Their broader argument is that polarization coverage can change what citizens think the political environment is like, which then shapes how they interpret disagreement. Under those conditions a citizen who would otherwise treat opponents as reasonable can come to treat them as alien, dangerous, or simply unreachable.

This is one reason the tracking problem is not solved by insisting that citizens only need to track constitutional essentials. If the public environment trains citizens to interpret disagreement as tribal conflict, then the civic stance public reason requires becomes harder to sustain. Levendusky and Malhotra show that perceptions of polarization affect affective attitudes in ways that are not tightly tethered to policy substance \parencite[p.~290]{levenduskyMalhotra2016polarization}. That is a direct threat to the civic trust that a public-justification practice requires.

Group polarization dynamics further amplify drift. Sunstein argues that deliberating groups tend to move toward more extreme positions in the direction indicated by their members' pre-deliberation tendencies \parencite[p.~176]{sunstein2002groupPolarization}. Cross-cutting exposure can counteract this effect. Mutz shows that cross-cutting networks can increase tolerance and political legitimacy while also changing the incentives for political participation \parencite[p.~111]{mutz2002crosscutting}. She also shows that exposure to disagreement can reshape how citizens evaluate political choices \parencite[p.~122]{mutz2002crosscutting}. In a high-sorting environment where cross-cutting exposure declines, group polarization becomes a likely pattern. Citizens can remain sincere and normatively committed to fair terms while inhabiting interpretive worlds that train their attention toward different features of the political landscape. The tracking task becomes harder because citizens do not merely disagree about values. They disagree about which facts and which considerations are salient.

These mechanisms interact. Cascades can shift what citizens think is empirically true about policies. Polarization can shift which considerations citizens treat as morally decisive. The combined effect is drift in the public coordinates citizens would need to orient themselves within a shared justificatory practice.

\section{Cognitive Constraints and Justificatory Drift}

The previous section concerned the public sphere. This section concerns a deeper limitation: the kind of cognition Rawlsian stability presupposes, and the kind of cognition human beings actually have. Here the internal-critique framing matters. I am not using cognitive science to redefine the reasonable. I am using it to ask whether the epistemic role stability-through-justification assigns to reasonable citizens is feasible given the way human beings process information under constraint.

Predictive-processing approaches begin from a simple idea: perception is not passive reception but active prediction. The brain continuously generates expectations about incoming signals and updates its model when those expectations are violated. Clark argues that this architecture treats perception as a process of controlling prediction error rather than constructing a complete mirror of the world \parencite[p.~190]{clark2013whatever}. Friston develops this point through a hierarchical account of cortical processing, emphasizing top-down prediction and bottom-up error correction \parencite[p.~816]{friston2005theory}. The details are not needed here. The upshot is that cognition is a tradeoff. It compresses. It discards information that is too costly to maintain. It is built to support workable action in a noisy world, not to provide a high-fidelity map.

Compression is not a defect. It is a condition of tractable intelligence. Ganguli and Sompolinsky show how high-dimensional signals can be encoded in compressed form while preserving enough structure for prediction and inference \parencite[pp.~491--492]{ganguli2012compressed}. Ahmad and Hawkins describe sparse distributed representations as a mechanism that yields robustness and generalization by representing information in high-dimensional patterns with only a small fraction of active components \parencite[p.~1]{ahmadHawkins2016sdr}. In stable environments such representations support reliable recognition. In environments where the structure of the signal changes slowly across multiple dimensions, compression can become a liability. Small shifts can be treated as noise until they accumulate into a rupture.

This is where drift becomes a defeater. A citizen can remain locally coherent. She can maintain a justificatory picture in which her public reasons appear to support the institutions she endorses. Yet the world can drift beneath that picture. Changes in technology, institutional interpretation, and social meaning can alter what a policy actually does without producing a sharp prediction error for the individual. Friston emphasizes that prediction error is the signal that drives updating; when error is small or ambiguous, the system can preserve its prior model rather than revise it \parencite[p.~134]{friston2010freeenergy}. Political environments often generate precisely that kind of ambiguity. They supply enough noise and enough interpretive degrees of freedom that citizens can reinterpret new information to fit their priors rather than revising those priors.

Friston also stresses a tension between accuracy and complexity. A cognitive system can reduce prediction error by building an increasingly rich model of the world, but rich models have costs. They demand metabolic and computational resources. They also become fragile, because small changes can produce large cascades of re-interpretation. The practical result is that updating is selective. A system will often prefer a cheap repair to a deep revision, even when the deeper revision would be more accurate in the long run \parencite[p.~128]{friston2010freeenergy}. This matters for political justification because many of the changes that alter justificatory status are not sharp violations. They are slow shifts in how policies operate, how categories are drawn, and which consequences are salient. Selective updating is a sensible cognitive strategy in ordinary life. It becomes a vulnerability when legitimacy depends on citizens keeping their justificatory judgments aligned with a drifting world.


The pattern is familiar in political life. Citizens experience long periods in which the meaning of a principle seems stable, then a sudden crisis in which the same principle appears to license incompatible outcomes. When that happens citizens often do not revise their justificatory commitments with fine-grained care. They adopt global narratives that restore coherence. The details vary. The structure does not. A legitimacy route that depends on steady, citizen-facing tracking of justificatory status is vulnerable to this cognitive architecture.

Rational ignorance sharpens the point. Somin argues that voter ignorance is often rational: citizens have little incentive to acquire detailed political information when the probability of affecting outcomes is low \parencite[p.~77]{somin2013democracy}. Rational ignorance does not mean citizens are stupid or vicious. It means the epistemic environment does not reward high-resolution tracking for most individuals. This does not refute public reason. It explains why public reason's citizen-facing stability ambition has hidden costs. Even if citizens want to be civically responsible, the structure of democratic decision-making does not reward most forms of slow, careful epistemic maintenance.

The granularity reply returns here in its strongest form. Public reason concerns constitutional essentials. Citizens need not track policy details. They can rely on institutional division of labor. This blocks an overly demanding reading of tracking. It does not dissolve the problem. Even coarse-grained tracking requires citizens to identify whether institutions remain within the liberal family, whether courts and legislatures are acting on public reasons, and whether the basic structure continues to secure fair terms. Those are not microscopic tasks. They remain tasks that require stable public markers and some reliable sense of what institutional signals mean over time.

\section{Complexity and Political Drift}

So far I have emphasized failures. A complexity-aware account must also take stabilizers seriously. Complex adaptive systems do not only drift. They can exhibit negative feedback loops that dampen shocks and restore equilibrium patterns. The question is what kind of stabilization Rawlsian stability requires, and whether the stabilizers available in modern politics are suited to that task.

Complexity enters the story in three ways. Exogenous drift comes from outside the political system: technological change, economic shocks, demographic shifts, geopolitical events. Endogenous drift comes from within: strategic incentives, polarization dynamics, institutional adaptation, and the way policies change the conditions under which later policies are evaluated. Non-linear drift occurs when the system crosses thresholds and reorganizes quickly. Gaus argues that the open society generates autocatalytic diversity: freedoms multiply social possibilities, which in turn generate further divergence \parencite[p.~102]{gaus2021open}. That picture makes it hard to treat the justificatory landscape as stationary.

Institutional complexity also frustrates top-down control. Scott's warning is that state schemes to make society legible often flatten local knowledge and destroy the very practices that make social systems resilient \parencite[pp.~2--3]{scott1998seeing}. A Rawlsian might reply that public reason is not a planning scheme. Yet the legitimacy ambition can still be threatened by the same dynamic. If the mapping from reasons to outcomes is complex and context-sensitive, then making that mapping publicly legible is difficult. The attempt can collapse into slogans, proxies, and category errors.

Negative feedback loops matter here. Polycentric governance can create self-correcting mechanisms. Ostrom argues that complex governance problems are often better handled by polycentric arrangements: overlapping centers of decision-making that enable experimentation, learning, and local adaptation \parencite[pp.~642--643]{ostrom2010polycentric}. Page emphasizes that diversity can function as insurance: groups with diverse perspectives can outperform homogeneous groups under certain conditions because they explore a broader space of solutions \parencite[p.~123]{page2011diversity}. The same point appears elsewhere in Page's discussion of how diversity can expand problem-solving capacity when the search space is rugged \parencite[p.~161]{page2011diversity}. Gaus argues that an open society requires rule-based moral expectations that allow strangers to coordinate without sharing a single thick conception of virtue \parencite[p.~2]{gaus2016rulebased}. These are real stabilizers. They suggest that drift can be tamed, at least episodically.

The deeper thought is that public justification is itself a rule-governed practice. Gaus argues that in a diverse moral world we should expect persistent disagreement about ultimate values, so the task of justification is not to discover a uniquely correct ordering of values but to secure a social order that can be justified to diverse agents under shared constraints \parencite[pp.~42, 283]{gaus2011order}. That picture is friendly to the Rawlsian ambition, but it is also less sanguine about the stability of shared justificatory understandings. If justification depends on a complex ecology of social rules, then drift can occur not only in policy outcomes but in the public meaning of justificatory norms themselves. That is another way the landscape can dance.


The difficulty is that negative feedback can restore order without restoring public justification. A system can correct errors through institutional adaptation that is opaque to citizens, or through bargaining equilibria that citizens accept as inevitable. Rawls wants a stability that expresses citizens' endorsement of public reasons, not merely a system that avoids collapse. A polycentric system can be more resilient while also being harder to narrate in a single justificatory story. Diversity can improve problem-solving while also widening disagreement about what counts as success. Rule-based baselines can stabilize expectations while also becoming ritualized in ways that mask substantive drift. Stabilizers can slow the dance without making it publicly visible.

A complexity-aware Rawlsianism therefore requires more than optimism about self-correction. It needs design commitments. The promising direction is institutional diversification and experimentation, combined with feedback mechanisms that allow citizens to contest and revise justificatory claims. Holland emphasizes that complex adaptive systems depend on layered signal-processing and boundary mechanisms that allow local units to respond without centralized control \parencite[pp.~1--2]{holland2012signals}. A political analogue would treat public justification as an ongoing practice that requires institutional supports to detect drift, to prevent capture, and to keep justificatory standards publicly inspectable.

The appeal of this direction can be seen by contrast with simpler constitutional models. Buchanan and Tullock offer a model of constitutional choice that aims to secure agreement on decision rules under conditions of disagreement \parencite[pp.~6--7]{buchanan1962calculus}. The minimalism is part of its attraction. Yet it lacks Rawls's moral ambition: the aim is not stability for the right reasons, but mutually advantageous cooperation. The comparison is instructive. If one wants Rawls's moral aim without Rawls's cognitive optimism, one needs institutions that can make the justificatory landscape legible enough for citizens to orient themselves without requiring each citizen to privately model the whole system.

That is the point at which my argument turns constructive. Before moving there, Rawls must be given every chance to recover on his own terms.

\section{The Best Rawlsian Recovery and What Survives It}

A critic of Rawlsian stability can too easily win by caricature. The relevant question is how far Rawls's own resources go once the tracking and landscape-stability assumptions are made explicit.

Start with the well-ordered society. Rawls does not imagine public reason operating in an epistemic void. A well-ordered society includes institutions and a public political culture that support citizens' sense of justice. One might think that these features function as cognitive stabilizers. If citizens share a political conception, and if institutions reliably realize it, then the epistemic burden on individuals declines. Citizens can treat institutional signals as trustworthy indicators of justificatory status.

Add the division of justificatory labor. Public reason binds most tightly for officials. Citizens may rely on courts, legislatures, and public norms to do much of the justificatory work. This is not an evasion. It reflects Rawls's awareness that ordinary citizens cannot be asked to perform the same justificatory tasks as judges. Rawls connects public reason to the duty of civility, asking citizens to see themselves as ideal legislators and to hold officials accountable \parencite[p.~444]{rawls2005political}. That is already a division of labor: citizens participate as overseers and critics rather than as technical reasoners on every issue.

Add educative institutions. A stable liberal order is not merely a set of rules. It is a social world that shapes citizens' capacities and expectations. If civic education, public discourse norms, and institutional design reduce informational volatility, then the justificatory landscape becomes more stable, and tracking becomes easier.

These recovery routes address the most obvious worry: that my critique assumes citizens must track too much, too quickly, and at too fine a granularity. They are also consistent with Rawls's restriction of public reason to constitutional essentials. The best version of the reply insists that stability does not require citizens to monitor every empirical shift. It requires that they can recognize whether the constitutional order remains within liberal bounds.

Weithman gives this reply its most sophisticated formulation. When Rawls concedes that an overlapping consensus is likely to be achieved over a family of liberal conceptions, stability cannot mean agreement on a single fixed conception. It becomes a matter of remaining within a range of acceptable conceptions while maintaining citizens' willingness to comply for the right reasons \parencite[p.~198]{Weithman2023FixedPoints}. Weithman develops this as a set-valued equilibrium: stability can be understood as the persistence of a well-ordered society within a \enquote{ball} of liberal conceptions rather than as a fixed point \parencite[p.~206]{Weithman2023FixedPoints}. This move directly engages the granularity parameter. It reduces the tracking burden by allowing variation within acceptable bounds.

Rawls's debate with Habermas helps illuminate what this recovery is trying to preserve. Habermas distinguishes three normative models of democracy and treats deliberation as a central source of legitimacy \parencite[p.~3]{habermas1994three}. Rawls's reply to Habermas seeks to preserve a role for democratic public reason without collapsing legitimacy into a comprehensive theory of rational discourse \parencite[p.~134]{rawls1995reply}. One can read Weithman's set-valued stability move as part of the same effort: to preserve legitimacy as a public practice while acknowledging that pluralism prevents full convergence on a single justificatory point.

What does this rescue genuinely fix? It blocks a naive reading of public reason as requiring constant high-resolution justificatory updating by every citizen. It makes room for institutional mediation. It recognizes that liberal legitimacy can be maintained even when citizens endorse different liberal conceptions, so long as those conceptions fall within the appropriate family. It also acknowledges that stability is an equilibrium notion, which helps explain why Rawlsian legitimacy is fragile under certain shifts in expectations.

Yet something survives even this best-case reconstruction. Tracking and landscape stability relocate rather than disappear.

Citizens must still be able to recognize whether institutions remain within the liberal family. That task does not reduce to knowing the constitution's text. It requires judgment about institutional interpretation, about the practical meaning of rights under changing social conditions, and about whether officials are acting from public reasons. Rawls claims that public reason can be exemplified by the supreme court in a constitutional regime, and he treats that exemplification as part of the public legitimacy story \parencite[p.~133]{rawls1995reply}. That claim is intelligible only if citizens can, at least in a coarse way, track whether the court's constitutional reasoning remains within the liberal family. Otherwise the appeal to public reason becomes a professional self-description rather than a public legitimacy route.

Rawls himself seems aware of this dependence. In \emph{Public Reason Revisited} he describes public reason as a shared ideal of citizenship and insists that public officials should be prepared to offer public reasons for their political decisions \parencite[p.~768]{rawls1997public}. The point of doing so is not exhausted by the official's own conscience. It is supposed to sustain a public political culture in which citizens can recognize justificatory reasons when they are offered, and can criticize officials when those reasons are missing. The official's discipline therefore presupposes an audience that can receive and assess, at least at a coarse level, what is being offered.


The set-valued stability move also requires publicly intelligible markers of family membership. If stability is remaining within a ball of acceptable conceptions, citizens need a way of telling when the society has drifted outside that ball. Weithman notes that contemporary instability worries have returned because citizens no longer experience liberal democracy as secure \parencite[p.~1187]{Weithman2024Rescuing}. A set-valued equilibrium does not supply markers by itself. It shifts the work to a different level: instead of tracking a point, citizens track boundaries.

Institutional mediation can also produce its own landscape instability. A division of labor can reduce individual cognitive burdens, but it also creates incentives for strategic actors to shape the informational environment. When justificatory authority is concentrated in institutions, citizens must trust those institutions. Vallier argues that polarization corrodes trust in government and in fellow citizens, making it harder to sustain cooperative political life \parencite[pp.~1, 3]{vallier2020trust}. The Rawlsian recovery package depends on trust as much as on reason. Trust is itself vulnerable to mediated drift.

The most defensible conclusion is therefore conditional and calibrated. Rawlsian recovery is not impossible in principle. Under conditions where institutional signals remain legible, where civic education and public culture stabilize expectations, and where the time horizon of demanded stability is long enough for citizens to update, a set-valued stability account can make the legitimacy story plausible. Under conditions of rapid short-cycle drift, high mediation, and polarized affect, the same story becomes fragile. The core tracking problem does not vanish. It moves.

This is also where the contrast with post-Rawlsian convergence liberalism becomes useful. Vallier argues that public justification can be sustained through convergence: different citizens can endorse the same coercive law for different reasons \parencite[p.~191]{vallier2011diss}. Gaus argues that in an open society moral diversity is autocatalytic and that justificatory unity is therefore hard to sustain \parencite[p.~14]{gaus2021open}. These views loosen a familiar Rawlsian demand for shared reasons. Yet they inherit the tracking problem. If public justification is sustained by multiple justificatory routes, citizens still need public ways of recognizing when laws remain justified and when justificatory routes have been distorted. Convergence can reduce one sort of burden while increasing the navigational task.

\section{Institutions as Cognitive Prosthetics}

If the tracking burden cannot be eliminated, it can be redistributed. That is the constructive pivot of the thesis. Institutions can function as cognitive prosthetics. They can reduce the cognitive burden on individuals by providing reliable public signals, audits, and correction mechanisms that make drift more visible and more contestable. The aim is not to make institutions think for citizens. The aim is to make it possible for citizens to see and contest the justificatory status of coercion without requiring each person to privately maintain a detailed map of a drifting landscape.

Hart's distinction between primary and secondary rules is a useful analogy. Primary rules govern conduct. Secondary rules govern how rules are recognized, changed, and adjudicated. Hart argues that a mature legal system requires secondary rules to avoid uncertainty, staticity, and inefficiency \parencite[pp.~92--93]{hart2012concept}. A public-reason analogue would treat justificatory practices as requiring second-order institutional supports: procedures for public diagnosis, for revision, for contestation, and for making the criteria of justification visible.

Hart's point can be put in Rawlsian terms. If citizens are to see law as authoritative, they need not only first-order rules that tell them what to do, but also publicly recognizable procedures for identifying which rules are valid, how rules can be changed, and how disputes about meaning are settled \parencite[pp.~79--80]{hart2012concept}. A public reason regime adds a further demand. Citizens must also be able to tell whether coercion is being exercised under publicly justifiable reasons. Without institutional supports for that recognition, the duty of civility becomes aspirational rhetoric rather than a lived political practice.


Polycentric design strengthens the case. Ostrom's analysis of polycentric governance emphasizes overlapping, independent centers of decision-making that allow for learning and error correction \parencite[p.~643]{ostrom2010polycentric}. A polycentric epistemic scaffold would not be a single ministry of public reason. It would be a set of institutions with partly overlapping functions: independent review bodies, transparency requirements, structured venues for contestation, and diverse sources of public information. The point is not redundancy for its own sake. Redundancy reduces capture risk and increases the chance that drift is detected before it becomes catastrophic.

Deliberative democrats have long argued that institutions can improve epistemic performance by structuring discussion and accountability. Bohman, reflecting on the \enquote{coming of age} of deliberative democracy, emphasizes that the theory has increasingly taken seriously the institutional conditions under which deliberation can function in real political settings \parencite[pp.~400--401]{bohman1998comingage}. Lafont adds a constraint that is directly relevant here: democratic legitimacy cannot be secured by shortcuts that ask citizens to defer to outcomes they cannot contest, even when those outcomes are produced by well-designed deliberative mini-publics \parencite[p.~3]{lafont2020democracy}. A cognitive-prosthetic approach should be read as aligning with Lafont's concern. The scaffolding must preserve the citizen's standing as an agent who can contest, revise, and authorize, not merely as a beneficiary of experts' epistemic competence.

A public-reason prosthetic also has to take account of complexity. Page's argument about diversity as insurance suggests that institutions should preserve heterogeneity rather than suppress it, because heterogeneity is part of what allows a system to search a rugged problem space \parencite[p.~130]{page2011diversity}. Gaus's account of rule-based order suggests that strangers can cooperate under shared norms without sharing a thick conception of virtue \parencite[p.~3]{gaus2016rulebased}. Holland's analysis of signals and boundaries suggests that complex systems work by layering feedback and by preventing any single part from dominating the whole \parencite[p.~2]{holland2012signals}. Together these ideas point toward a design ethos: public justification should be supported by multiple, contestable, publicly visible mechanisms rather than by a single centralized epistemic authority.

This proposal invites the paternalism objection. The objection expresses a core liberal worry. If institutions filter what reasons citizens encounter, do they not smuggle elite control into the very practice meant to express equal citizenship? Do they not replace citizens' judgment with curated narratives, thereby undermining the standing of citizens as co-authors of law?

The objection has bite because modern politics already contains forms of epistemic control. Scott warns that legibility schemes can become tools of domination: by simplifying what is measured, the state can impose categories that shape social reality \parencite[p.~309]{scott1998seeing}. A public-reason prosthetic could become an epistemic bureaucracy that decides what counts as a public reason and what does not. That would be a betrayal of the project.

The response must therefore be structural. A cognitive prosthetic is not paternalistic if it satisfies four conditions.

It must be transparent. The criteria by which reasons are assessed must be public, and the institution must explain how it reaches its judgments. Transparency aims to expose reasons, not to hide them.

It must be contestable. Citizens and civil society must have structured avenues to challenge justificatory claims and to demand revision. Lafont stresses that democratic legitimacy requires that citizens not be expected to blindly defer to decisions made by others; contestation is part of what makes citizens authors rather than subjects \parencite[p.~3]{lafont2020democracy}.

It must be polycentric. No single institution should have a monopoly on justificatory filtering. Multiple independent bodies reduce the risk that a single elite can capture the scaffolding. Ostrom's polycentricity supplies the core anti-monopoly model \parencite[pp.~644--645]{ostrom2010polycentric}.

It must be constitutionally authorized and revisable. Citizens must be able to endorse, revise, and if necessary replace these scaffolds through ordinary democratic mechanisms. Without authorization and revisability, the prosthetic becomes a technocratic imposition rather than a public practice.

These safeguards do not eliminate disagreement. They aim to make drift visible and politically manageable. They shift public reason from a demand on individual cognition to a property of institutional design. A society is more capable of stability for the right reasons when it has public, contestable, polycentric mechanisms that help citizens see whether coercion remains publicly justifiable.

Landemore's work on collective intelligence supports the general direction. Democratic decision-making can outperform small elites when institutions harness diversity of information and perspective \parencite[pp.~3--4]{landemore2013democratic}. The challenge is that diversity can also generate noise and polarization. A prosthetic design aims to preserve the epistemic benefits of diversity while providing structures that prevent cascades and affective polarization from dominating the justificatory environment.

\section{Conclusion}

The thesis began with a Rawlsian ambition: legitimacy through stability for the right reasons. I have tried to show why that ambition has invisible feasibility conditions. Citizens must be able to track justificatory status over time, and the justificatory landscape must be stable enough for tracking to work. In modern mediated politics both conditions are often fragile. Empirical drift and affective polarization can defeat tracking even among sincere, normatively reasonable citizens. Cognitive compression can preserve coherence while allowing the world to drift beneath justificatory beliefs. Complexity can generate stabilizers, but stabilizers can restore order without restoring public justification.

None of this proves that Rawls is false, and it does not warrant cynicism about deliberation. Estlund is right to warn against a timid political philosophy that treats aspirational ideals as illegitimate because they are hard to realize \parencite[pp.~10--12]{estlund2019utopophobia}. The point here is narrower. If legitimacy is defined by stability-through-public-justification, then feasibility constraints delimit the domain in which that legitimacy route can plausibly function. Where the conditions fail, the theory must either restrict its ambition or redesign its institutional story.

The best Rawlsian recovery, especially in Weithman's set-valued stability terms, reduces the demand by allowing variation within liberal bounds. It does not erase the need for citizens to track boundary conditions and institutional legitimacy. That is why the constructive turn matters. Institutions can carry part of the justificatory load in a democratically acceptable way if they function as cognitive prosthetics rather than as epistemic controllers. The safeguards I have defended are not rhetorical decorations. They are the difference between scaffolding and paternalism.

Weithman argues that contemporary democratic instability should be treated as a genuine philosophical problem rather than as mere moral panic \parencite[p.~1187]{Weithman2024Rescuing}. One reason is that the failure mode can be quiet. A constitutional order can retain formal liberties while losing the public bases of justification that make those liberties feel secure. Rawls anticipates the affective version of this problem when he warns that citizens who no longer see the point of the public reason ideal can fall into resentment \parencite[p.~485]{rawls2005political}. The deeper worry is that citizens can continue to talk the language of public reason while the practical mapping from that language to coercive outcomes becomes opaque.

For Rawlsians the upshot is not a refutation of public reason but a forced choice about what exactly is being defended. If stability for the right reasons is meant to hold over long horizons, then institutions must do more epistemic work than Rawls sometimes acknowledges. If stability is meant to be a thinner virtue, holding only episodically in favorable periods, then the legitimacy ambition must be trimmed accordingly. The middle path is to accept stability as an equilibrium ideal, in Weithman's sense, while designing public, contestable mechanisms that make the relevant equilibrium conditions visible enough for citizens to sustain them.

For convergence liberals the lesson is parallel. Allowing diverse reasons to support the same law relaxes a demand for shared justificatory currency, but it increases the need for public methods of diagnosing drift and maintaining trust across divergent justificatory routes \parencite[p.~297]{gaus2011order}. In a plural order, the problem is rarely that citizens have no reasons. The problem is that the reasons become unmoored from what institutions are actually doing, and that unmooring is hard to detect from within any single justificatory perspective.

If stability through public reason is to be more than episodic, liberal democracies will need to treat the epistemic dimension of legitimacy as a design problem. Not as an invitation to replace normative theory with psychology, but as a recognition that public justification is a practice with cognitive and institutional preconditions. A liberal order that cannot help its citizens see what it is doing will not remain legitimate for long, even if it survives.

\printbibliography

\end{document}
